\chapter{绪\quad 论}
\label{chap:intro}

\section{研究背景与意义}
\subsection{研究背景}
近年来，深度学习与大语言模型在自然语言处理与科研辅助领域取得了显著进展\cite{vaswani2017attention,devlin2019bert}。在研究生撰写硕士或博士学位论文的实际场景中，往往需要兼顾 LaTeX 的公式与排版优势以及导师审阅时习惯使用的 Word 文档格式。

\subsection{研究意义}
构建统一、可追溯、支持多文件分章写作与高质量 Word 导出的学位论文写作平台，能够显著降低格式排版负担，提升师生协作效率与学术合规性。

\section{国内外研究现状及分析}
现有在线写作工具主要分为纯 LaTeX 编辑器与富文本在线文档两类\cite{he2016deep}，两者在公式排版、参考文献规范与导师批注体验上各有所长，但在双向工作流协同方面仍有明显断层。

\section{本文主要研究内容}
本文围绕多文件大论文写作与原生 Word 导出两大核心需求展开研究，重点解决跨文件导航、按章编号一致性以及复杂公式转换等关键问题。

\section{本文章节安排}
本文共分为五章：第~\ref{chap:intro}~章为绪论；第~\ref{chap:theory}~章介绍相关理论与技术；第~\ref{chap:method}~章阐述核心方法与实验分析；第~\ref{chap:system}~章介绍原型系统设计与实现；第~\ref{chap:conclusion}~章总结全文并展望未来工作。

\section{本章小结}
本章梳理了研究背景与意义，明确了本文的研究目标与整体章节结构。
